\section*{अध्याय \ref{ch:b2:biomolecules} --- जैव-अणु: ऐमीनो अम्ल, शर्कराएँ, लिपिड और न्यूक्लियोटाइड}

\begin{solution}{exo:b2:biomolecules:1}
\begin{itemize}
  \item pH 3 और 9 के बीच \omterm{def:b2:biomolecules:amino-acid}{उभयाविष्ट आयन} \ce{H3N+-CH(CH3)-COO-}।
  \item pH 1 पर धनायन \ce{H3N+-CH(CH3)-COOH}।
  \item pH 12 पर ऋणायन \ce{H2N-CH(CH3)-COO-}।
\end{itemize}
\end{solution}

\begin{solution}{exo:b2:biomolecules:2}
ग्लाइसीन: $(2.35 + 9.78)/2 \approx 6.1$। ऐलानीन: $(2.35 + 9.87)/2 \approx 6.1$।
% ledger: pka:glycine1, pka:glycine2, pka:alanine1, pka:alanine2
\end{solution}

\begin{solution}{exo:b2:biomolecules:3}\emergencystretch=2em
\ce{H2N-CH(CH3)-CO-NH-CH2-COOH}: ऐलानीन N-अंत्य प्रदान करता है, ग्लाइसीन C-अंत्य। Gly--Ala,
\ce{H2N-CH2-CO-NH-CH(CH3)-COOH}, एक संरचनात्मक समावयवी है: मुक्त ऐमीनो समूह ग्लाइसीन पर है।
\end{solution}

\begin{solution}{exo:b2:biomolecules:4}
\ce{NH2} बाईं ओर: L। प्राथमिकताएँ \ce{NH2} > \ce{COOH} > \ce{CH2OH} > \ce{H}: L-सेरीन
(\textit{S}) है।
\end{solution}

\begin{solution}{exo:b2:biomolecules:5}
समूह: N-अंत्य अमोनियम (लगभग 9), लाइसीन \omterm{def:b2:biomolecules:amino-acid}{पार्श्व शृंखला} (10.7), ऐस्पार्टेट \omterm{def:b2:biomolecules:amino-acid}{पार्श्व शृंखला} (3.9),
C-अंत्य कार्बोक्सिलिक समूह (लगभग 2)। pH 1: $+1 +1 + 0 + 0 = +2$। pH 7: $+1 +1 -1 -1 = 0$। pH 12:
$0 + 0 - 1 - 1 = -2$ (लाइसीन \omterm{def:b2:biomolecules:amino-acid}{पार्श्व शृंखला}, pH 12 से 1.3 इकाई नीचे, अधिकांशतः उदासीन है)।
\end{solution}

\begin{solution}{exo:b2:biomolecules:6}\emergencystretch=3em
थायोनिल क्लोराइड और ताप अन्य समूहों पर भी आक्रमण करते और, सबसे बढ़कर, किसी ऐमीनो अम्ल का अत्यंत अभिक्रियाशील \omterm{def:b2:acyl-substitution:derivative}{ऐसिल क्लोराइड}
आसानी से अपना विन्यास खो देता है ($\alpha$ हाइड्रोजन अम्लीय हो जाता है): \omterm{def:b2:biomolecules:peptide}{पेप्टाइड}
आंशिक रूप से रेसेमिक हो जाता। \omterm{def:b2:biomolecules:peptide}{युग्मन अभिकर्मक} कमरे के ताप पर, मृदु परिस्थितियों में, अम्ल को सक्रिय करता है।
\end{solution}

\begin{solution}{exo:b2:biomolecules:7}
O5 C2 से आबंधित: C2, C3, C4, C5 और O का वलय, पाँच परमाणु (एक फ़्यूरैनोस)। \omterm{def:b2:biomolecules:anomer}{हॉवर्थ प्रक्षेप} में वलय O
पीछे, C2 दाईं ओर, \ce{CH2OH} (C1) नीचे और \ce{OH} ऊपर ($\beta$); C3 पर \ce{OH} ऊपर (फ़िशर में
बाईं ओर), C4 पर \ce{OH} नीचे, C5 \ce{CH2OH} (C6) को ऊपर रखता हुआ।
\end{solution}

\begin{solution}{exo:b2:biomolecules:8}
$x_\alpha = (80.2 - 52.8)/(150.7 - 52.8) = 27.4/97.9 \approx 0.28$: 28~\% $\alpha$, 72~\% $\beta$।
\end{solution}

\begin{solution}{exo:b2:biomolecules:9}
माल्टोस अपने दूसरे ग्लूकोस पर एक मुक्त हेमीऐसीटैल बनाए रखता है, जो खुलकर ऐल्डिहाइड बनता है: अपचायी। सुक्रोस में
दोनों \omterm{def:b2:biomolecules:anomer}{ऐनोमरी कार्बन} \omterm{def:b2:biomolecules:glycoside}{ग्लाइकोसाइडी आबंध} में लगे हैं: कोई हेमीऐसीटैल नहीं, अपचायी नहीं। तनु अम्ल
\omterm{def:b2:biomolecules:glycoside}{ग्लाइकोसाइडी आबंध}, एक ऐसीटैल, का जल-अपघटन करता है: ग्लूकोस और फ़्रक्टोस मुक्त होते हैं, और दोनों अपचयन करते हैं।
\end{solution}

\begin{solution}{exo:b2:biomolecules:10}
ट्राइओलीन, \qty{885.45}{g/mol}, तीन \ce{KOH} (\qty{56.105}{g/mol}) का उपभोग करता है: $3 \times 56.105/885.45 =
0.190$~g प्रति g, अर्थात् \omterm{def:b2:acyl-substitution:saponification}{साबुनीकरण} मान 190। छोटी शृंखलाओं का अर्थ है उन्हीं तीन एस्टर समूहों के लिए
छोटा मोलर द्रव्यमान: प्रति ग्राम वसा के अधिक मोल, अधिक KOH।
% ledger: aw:C, aw:H, aw:O, aw:K
\end{solution}

\begin{solution}{exo:b2:biomolecules:11}
ग्लाइसीन के ऐमीन का \textit{tert}-ब्यूटॉक्सीकार्बोनिल कार्बामेट के रूप में, और ऐलानीन के अम्ल का
मेथिल एस्टर के रूप में रक्षण कीजिए। किसी कार्बोडाइइमाइड से युग्मन कीजिए। कार्बामेट को अम्ल से और एस्टर को तनु क्षारक से हटाइए
(लांबिक परिस्थितियाँ, अतः किसी भी सिरे को अकेले मुक्त किया जा सकता है): Gly--Ala।
\end{solution}

\begin{solution}{exo:b2:biomolecules:12}
5$'$-ACGCAT-3$'$ (प्रतिसमांतर पढ़ा गया: A, T से युग्मित होता है, G, C से)। युग्म: A--T, T--A, G--C, C--G, G--C, T--A:
$3 \times 2 + 3 \times 3 = 15$ हाइड्रोजन आबंध।
\end{solution}

\begin{solution}{pb:b2:biomolecules:1}\emergencystretch=8em
\textbf{1.} L-ऐस्पार्टिक अम्ल, L-फ़ेनिलऐलानीन और मेथेनॉल।
\textbf{2.} दोनों में \ce{COOH} शीर्ष पर, \omterm{def:b2:biomolecules:amino-acid}{पार्श्व शृंखला} नीचे (\ce{CH2COOH}, \ce{CH2C6H5}), \ce{NH2}
बाईं ओर।
\textbf{3.} दोनों के लिए (\textit{S})।
\textbf{4.} दो त्रिविमजनक केंद्र: चार त्रिविम समावयवी।
\textbf{5.} \ce{H2N-CH(CH2COOH)-CO-NH-CH(CH2C6H5)-COOCH3}: \omterm{def:b2:biomolecules:peptide}{पेप्टाइड आबंध} दोनों
अवशेषों को जोड़ता है; एस्टर C-अंत्य पर स्थित \ce{COOCH3} है।
\textbf{6.} स्वाद ग्राही किरैल हैं: वे एक त्रिविम समावयवी को बाँधते हैं, उसके दर्पण प्रतिबिंब या
अप्रतिबिंबी त्रिविम समावयवियों को नहीं, जैसा आरंभ में कहा गया था।
\textbf{7.} N-अंत्य \ce{NH3+} और ऐस्पार्टिक अम्ल का पार्श्व-शृंखला \ce{COOH}; C-अंत्य अम्ल
एक मेथिल एस्टर है, जिसमें कोई अम्लीय हाइड्रोजन नहीं।
\textbf{8.} pH 2 पर: \ce{NH3+} ($+1$), पार्श्व-शृंखला \ce{COOH} (0): $+1$।
\textbf{9.} pH 7 पर: $+1 - 1 = 0$।
\textbf{10.} pH 12 पर: $0 - 1 = -1$।
\textbf{11.} बदलने वाले समूहों के दोनों $\mathrm pK_a$ के बीच: $(3.9 + 9.9)/2 \approx 6.9$।
\textbf{12.} डाइपेप्टाइड में ऐमीन के कार्बन पर कार्बोक्सिलेट के बजाय एक ऐमाइड है: पड़ोसी
समूह और आवेश बदल जाते हैं, और $\mathrm pK_a$ भी (किसी \omterm{def:b2:biomolecules:peptide}{पेप्टाइड} का N-अंत्य अमोनियम
मुक्त ऐमीनो अम्ल के अमोनियम से स्पष्टतः अधिक अम्लीय है)। मॉडल आकृति देता है, यथातथ
मान नहीं।
\textbf{13.} अम्ल और ऐमीन पहले लवण बनाते हैं; गर्म करने से मिश्रण मिलते (Asp--Asp, लंबी शृंखलाएँ,
ग़लत कार्बोक्सिलिक समूह) और त्रिविमजनक केंद्र रेसेमिक हो जाते।
\textbf{14.} इसका ऐमीनो समूह और इसका पार्श्व-शृंखला कार्बोक्सिलिक समूह।
\textbf{15.} यह अम्ल के \ce{OH} को एक अच्छे निर्गामी समूह में बदल देता है, एक सक्रियित व्युत्पन्न जिस पर
ऐमीन कमरे के ताप पर आक्रमण करता है (ऐसिल प्रतिस्थापन)।
\textbf{16.} $\beta$-डाइपेप्टाइड, जो \omterm{def:b2:biomolecules:amino-acid}{पार्श्व शृंखला} से जुड़ा है: ऐस्पार्टेम का एक समावयवी, ऐस्पार्टेम नहीं।
\textbf{17.} ऐमीन का बेन्ज़िलॉक्सीकार्बोनिल कार्बामेट के रूप में और \omterm{def:b2:biomolecules:amino-acid}{पार्श्व शृंखला} का बेन्ज़िल एस्टर के रूप में रक्षण: दोनों
पैलेडियम पर हाइड्रोजन-अपघटन से एक साथ हटते हैं, जो मेथिल एस्टर को अछूता छोड़ता है (उसके
लांबिक)।
\textbf{18.} सक्रियित अम्ल का $\alpha$ हाइड्रोजन अम्लीय होता है; क्षारक उसे हटा देता है और त्रिविमजनक केंद्र
रेसेमिक हो जाता है।
\textbf{19.} ऐस्पार्टिक अम्ल का रक्षण (N पर कार्बामेट, \omterm{def:b2:biomolecules:amino-acid}{पार्श्व शृंखला} पर बेन्ज़िल एस्टर); मुक्त
$\alpha$-\ce{COOH} को \omterm{def:b2:biomolecules:peptide}{युग्मन अभिकर्मक} से सक्रिय कीजिए; L-फ़ेनिलऐलानीन का मेथिल एस्टर मिलाइए; हाइड्रोजन-अपघटन कीजिए।
\textbf{20.} मेथिल एस्टर और \omterm{def:b2:biomolecules:peptide}{पेप्टाइड आबंध}; एस्टर पहले, क्योंकि ऐमाइड कहीं कम अभिक्रियाशील है।
\textbf{21.} डाइपेप्टाइड Asp--Phe (मुक्त अम्ल) और मेथेनॉल।
\textbf{22.} मीठा अणु उपभुक्त हो जाता है: उसके जल-अपघटन उत्पाद उसका स्थान नहीं लेते।
\textbf{23.} $14 \times 12.011 + 18 \times 1.008 + 2 \times 14.007 + 5 \times 15.999 =
\qty{294.307}{g/mol}$।
\textbf{24.} $0.180/294.307 = \qty{6.116e-4}{mol}$, अर्थात् \qty{0.612}{mmol}।
\textbf{25.} प्रति ऐस्पार्टेम एक मेथेनॉल: $0.6116 \times 32.042 = \boldsymbol{\approx \qty{19.6}{mg}}$
मेथेनॉल।
\end{solution}
